\begin{exercise}[subtitle={Monotone functions have countable discontinuities \Coffeecup\Coffeecup\Coffeecup}]
    \label{ex: monotone functions have countably many discontinuities}
    Let \(F\colon \real \to \real\) be a non-decreasing function.
    \begin{enumerate}[label=(\alph*), noitemsep]
        \item\label{it: (monotone functions) left and right limits}
        Let \(x\in \real\). Prove that the limit
        \[
            F_-(x) \coloneq \lim_{n\to \infty} F(x_n)
        \]
        exists for any sequence \((x_n)_{n\in \nat} \subseteq \real\) with \(x_n < x\)
        and \(x_n \to x\). Prove that this limit is the same for any such
        sequence such that \(F_-(x)\) is well-defined.
        Also prove that \(F_+(x) \coloneq \lim_{y\downarrow x} F(y)\) is
        well-defined.

        \textbf{Hint:} Consider monotonically increasing sequences \(x_n\) first.
        Then prove that the liminf and limsup coincide for general sequences
        using monotone subsequences.

        \item\label{it: F continuous if and only if all coincide}
        Prove that \(F\) is continuous at \(x\) if and only
        if \(F(x) = F_-(x) = F_+(x)\). That is
        \(F\) only has jump discontinuities.

        \item\label{it: (monotone functions) countable jumps, compact interval}
        Prove that \(F\) restricted to a compact interval \([a,b]\)
        has only countably many jumps (places where \(F_-(x) < F_+(x)\)).

        \textbf{Hint:} How many jumps are there of size greater than \(1/n\)?

        \item\label{it: (monotone functions) countable jumps, real line}
        Prove that \(F\) has at most countably many discontinuities.
    \end{enumerate}
\end{exercise}

\begin{solution}
\begin{enumerate}[wide=0pt]
    \item[\ref{it: (monotone functions) left and right limits}]
    Let \(x_n\uparrow x\) be a monotonically increasing sequence. Then
    \(F(x_n)\) is also monotonically increasing and thus assumes a limit.
    We will prove \(\lim_n F(x_n) \le \lim_n F(y_n)\)
    for another monotonically increasing sequence \(y_n\uparrow x\).
    Since \(x_n\) and \(y_n\) are exchangeable we then also have equality.
    Now for any \(n\in \nat\) there exists an \(m_n\in \nat\) such that
    \(x_n \le y_{m_n} \le x\) as \(x_n < x\) and \(y_n\to x\). Consequently
    since \(F\) is non-decreasing we have
    \[
        \lim_{n\to\infty} F(x_n) \le \lim_{n\to\infty}F(y_{m_n}) \overset{\text{subsequence}}= \lim_{n\to\infty}F(y_n).
    \]
    Thus we have proven that the limit of monotonically increasing
    sequences is unique and denoted by \(F_-(x)\). Let \(x_n \to x\) with \(x_n < x\) be an
    arbitrary sequence. Let \(x_{n_k}\) be a subsequence with
    \[
        F(x_{n_k}) \to \liminf_{n\to\infty} F(x_n)
    \]
    Since \(x_{n_k} \to x\) and \(x_{n_k} < x\) we can find a monotonically increasing subsequence.
    This subsequence converges to \(F_-(x)\) and since it must have the
    same limit as the original sequence \(x_{n_k}\) we have
    \[
        \liminf_{n\to\infty} F(x_n) = F_-(x).       
    \]
    The same argument works for the limsup. Thus the limit exists and is
    given by \(F_-(x)\). The exact same argument also works for
    \(F_+(x)\) with decreasing sequences.

    \item[\ref{it: F continuous if and only if all coincide}]
    Clearly we have \(F_-(x) = F_+(x)\) if \(F\) is continuous.  So we
    only need to prove \(F\) continuous in \(x\) when this equality
    holds true.  Note that
    \(F_-(x) \le F(x) \le F_+(x)\) since \(F\) is non-decreasing.
    So if we have the equality then by a sandwich argument they are
    both equal to \(F(x)\). Let us assume equality. To prove
    sequential continuity in \(x\) we choose a sequence \(x_n\to x\).
    This sequence can be split into the parts that are smaller than \(x\)
    converging to \(F_-(x)\), the parts that are equal to \(x\),
    converging to \(F(x)\) and the parts that are larger than \(x\)
    converging to \(F_+(x)\). Since all these limits coincide the entire
    sequence converges to \(F(x)\). Consequently \(F\) is continuous in
    \(x\).

    \item[\ref{it: (monotone functions) countable jumps, compact interval}]
    Since \(F(b) - F(a)\) is finite and \(F\) monotonically increasing,
    there are at most \(\lceil n(F(b)-F(a))\rceil\) jumps of size
    \(\frac1n\). Thus for any \(n\in \nat\) the number of jumps of greater
    size is finite. This implies that the discontinuity set
    \[
        D_F = \set{x: F_-(x) \neq F_+(x)} = \bigcup_{n\in \nat} \set[\big]{x: F_+(x) - F_-(x) > \tfrac1n}
    \]
    is countable.

    \item[\ref{it: (monotone functions) countable jumps, real line}]
    Since \(F\) only has countably many discontinuities in each
    interval \([n, n+1]\) for \(n\in \integer\) and the countable union
    of countable sets is countable, we have that the set of all
    discontinuities is at most countable.
\end{enumerate}
\end{solution}
